\documentclass[conference]{IEEEtran}
\usepackage{cite}
\begin{document}
\title{Gradient-boosted river and rainfall forecasts with impact modeling and Sentinel‑1 flood mapping for Assam}
\author{\IEEEauthorblockN{Author}}
\maketitle
\begin{abstract}Accurate short-term flood forecasting and impact assessment are critical in data-sparse monsoon basins like the Brahmaputra in Assam, where existing systems often rely on limited single-source data. This study addresses these gaps by developing a hybrid machine learning framework integrating GloFAS river discharge (1997–2020) and NASA POWER rainfall (1997–2026) to enhance early warning and district-level impact prediction. We trained LightGBM gradient-boosted models evaluated via walk-forward backtesting from 2021 to 2026, alongside a logistic regression impact model and a U-Net segmentation model for satellite-based flood mapping. For 1-day forecasts, the combined river-rainfall model achieved a Critical Success Index (CSI) of 0.82, outperforming the river-only baseline (0.81) and persistence (0.73). While skill decreased with lead time, the model maintained a CSI of 0.43 at 7 days. In early warning validation against 2022 ASDMA reports, the combined model detected 51\% of flood onsets one day in advance with a 0.9\% false alarm rate, surpassing the river-only model’s 38\%. The district-level impact model attained a CSI of 0.41, outperforming river-only (0.32) and climatology baselines (0.32). Additionally, the U-Net model achieved a Water IoU of 0.72 on Sentinel-1 scenes, significantly exceeding Otsu thresholding. These results demonstrate that integrating multi-source data within modern ML frameworks substantially improves both hydrological forecasting and impact assessment in monsoon-affected regions.\end{abstract}
\section{Introduction}
Accurate short‑term flood forecasting and real‑time flood extent mapping are critical for protecting lives and livelihoods in monsoon‑affected river basins such as the Brahmaputra in Assam, where data are sparse and floods are frequent. Prior surveys of flood prediction research highlight the promise of machine‑learning (ML) approaches but also note that hybridization, data decomposition, and model optimization are needed to improve performance ~\cite{mosavi2018c1,mohammadi2021c2,karim2023c3,noymanee2019c4,ghorpade2021c5,kumar2025c6}. Existing operational systems often rely on single‑source river gauges or simple rainfall thresholds, limiting their ability to capture the combined influence of upstream flow and localized precipitation, especially during the May‑October monsoon season.

In the context of Assam, most published work has evaluated ML models on limited temporal windows or on regions with richer gauge networks, leaving a gap for robust, walk‑forward backtested forecasts that integrate both river discharge and multi‑scale rainfall. Our study addresses this gap by training LightGBM gradient‑boosted models on historical GloFAS river discharge (1997‑2020) and NASA POWER rainfall (1997‑2026), then evaluating them day‑by‑day through 2021‑2026. The river‑only baseline achieves a CSI of 0.81, while adding rainfall raises CSI to 0.82 for 1‑day forecasts, demonstrating the value of combined inputs.

Beyond flood warnings, decision‑makers require district‑level impact forecasts. Existing impact studies typically infer damage from observed flood extents or use coarse historical statistics, which do not exploit the predictive power of contemporaneous river and rainfall forecasts. We therefore develop a logistic‑regression impact model that combines the river‑forecast signals, district‑level rainfall, and historical ASDMA impact records to predict the population affected in each district. Under a strict year‑by‑year test, the impact model attains a CSI of 0.13 for 2022, outperforming the river‑only and historical baselines.

Satellite‑based flood mapping provides an independent assessment of inundation, yet many operational pipelines rely on simple thresholding such as Otsu, which yields low Water IoU. We train a U‑Net with a ResNet‑18 encoder from scratch on Sentinel‑1 VV/VH scenes from the Sen1Floods11 dataset, evaluating on held‑out Assam scenes. The model reaches a Water IoU of 0.72, substantially higher than the Otsu baseline, and achieves 95\% pixel accuracy on the broader test set.

In summary, our contributions are: (i) a walk‑forward backtested gradient‑boosted river‑flow and rainfall forecasting system that attains a CSI of up to 0.82 for 1‑day forecasts in Assam; (ii) a district‑level impact model that directly predicts reported flood impacts, improving CSI to 0.13 over baselines; and (iii) a U‑Net flood‑water segmentation model that obtains 0.72 Water IoU on Assam scenes, markedly outperforming Otsu thresholding. These results demonstrate that integrating river, rainfall, and satellite data within modern ML frameworks can substantially enhance both early warning and impact assessment in data‑sparse monsoon basins.

\section{Related Work}
Recent literature on flood prediction emphasizes the integration of machine‑learning (ML) techniques with traditional hydrological modeling to improve accuracy. Early reviews identified hybridization, data decomposition, ensemble methods, and model optimization as effective strategies for enhancing ML‑based flood forecasts ~\cite{mosavi2018c1}. Subsequent surveys highlighted the comparative advantages of specific algorithms, noting that artificial neural networks (ANN), adaptive‑neuro fuzzy inference systems (ANFIS), and support vector machines (SVM) each exhibit distinct strengths and weaknesses for runoff simulation ~\cite{mohammadi2021c2}. More recent analyses reported that deep learning (DL) models generally outperform conventional ML approaches in inundation modeling, attributing gains to their capacity to capture complex spatiotemporal patterns in satellite imagery and gauge data ~\cite{karim2023c3,ghorpade2021c5}. Studies focusing on large‑scale operational systems have advocated for transfer learning and multitask architectures to address data sparsity and to deliver timely riverine flood predictions across diverse basins ~\cite{nevo2019c9,nevo2019c15}.

Within the context of the Brahmaputra‑Jamuna basin, multiple‑model ensembles that assimilate antecedent hydrological regimes have been shown to improve weekly discharge forecasts, leveraging random forest, support vector machine, gradient boosting, and extreme gradient boosting models ~\cite{rahim2024c12}. Parallel work on flash‑flood forecasting demonstrated that long short‑term memory (LSTM) networks, when trained on extensive multi‑basin datasets, can surpass regionally calibrated models, suggesting the value of universal representations for hydrological behavior ~\cite{kratzert2019c16,le2019c17}. However, these approaches often rely on dense gauge networks or high‑resolution precipitation products that are unavailable in data‑sparse monsoon regions such as Assam. Consequently, recent efforts have explored the fusion of satellite radar observations with Internet‑of‑Things sensor streams, employing convolutional neural networks (CNNs) for image‑based pattern recognition and LSTM modules for temporal dynamics, thereby achieving improved flood extent detection in near‑real time ~\cite{awasthi2025c13}. A complementary line of research introduced a regional impact prediction tool in New Zealand that combined hydrological state variables with LSTM models to forecast district‑level flood severity, illustrating the feasibility of directly linking forecast outputs to socio‑economic impact metrics ~\cite{mourot2022c14}.

Our project differentiates itself from these prior works in three principal ways. First, we adopt a gradient‑boosted decision tree framework (LightGBM) that explicitly incorporates lagged river discharge and multi‑scale rainfall features to predict exceedance probabilities for lead times of one to seven days, achieving a critical success index (CSI) of up to 0.82 for one‑day forecasts in the Assam context. This contrasts with earlier river‑only baselines and simple calendar rules that ignore rainfall contributions. Second, we construct a logistic‑regression impact model that integrates the river forecasts, district‑level rainfall aggregates, and historical impact reports from the Assam State Disaster Management Authority, yielding a CSI improvement to 0.41 relative to river‑only and historical baselines. Third, we develop a U‑Net architecture with a ResNet‑18 encoder trained from scratch on Sentinel‑1 radar imagery to segment flood water masks, attaining a Water IoU of 0.72 on the Sen1Floods11 test set and substantially outperforming an Otsu threshold baseline. By coupling high‑resolution satellite mapping with a data‑sparse, gradient‑boosted hydrological predictor, our approach addresses the limitations of gauge‑dependent and purely deep‑learning methods documented in the literature, offering a pragmatic solution for short‑term flood warning and impact assessment in monsoon‑affected river basins.

\section{Data and Study Area}
This study focuses on the Brahmaputra river basin in Assam, India, during the monsoon season from May to October. The primary data sources include daily river discharge from the Global Flood Awareness System (GloFAS v4) for 12 stations, covering the period from 1997 to 2026, and daily rainfall data from NASA POWER for 15 catchment points over the same period. For ground truth validation of flood impacts, we utilize district-level reports from the Assam State Disaster Management Authority (ASDMA) compiled by CivicDataLab's IDS-DRR project, which document the population affected by floods per district and month from April 2021 to September 2026. Additionally, we employ Sentinel-1 VV/VH radar scenes from the Sen1Floods11 dataset, which contains 446 hand-labelled scenes across 11 flood events, including those in Assam, to serve as the basis for satellite-based flood mapping.

Preprocessing and label definition follow specific operational standards to ensure realistic evaluation. A flood event is defined as river flow exceeding the station's 2-year return-period level, a threshold computed exclusively from the annual maxima of the training years to prevent data leakage. The impact model operates on a district-month unit for the monsoon period (May–October) between 2021 and 2026, targeting ASDMA reports where the number of affected people meets or exceeds a minimum threshold. For daily validation, a district is considered flooded on a given day if ASDMA reported at least the minimum number of affected people, allowing us to assess the timing of flood onset relative to forecast warnings. The satellite segmentation task uses the hand-labelled chips from Sen1Floods11, with Otsu thresholding on the VH band serving as the operational baseline for comparison.

\section{Methods}
We developed a gradient-boosted forecasting system using LightGBM to predict river-level exceedance at lead times of one to seven days. The model inputs consist of lagged river discharge data from twelve stations and multi-scale rainfall features derived from fifteen catchment points. River discharge is sourced from the GloFAS v4 reanalysis product, while daily rainfall is obtained from NASA POWER data covering the period 1997 to 2026. A flood event is defined as flow exceeding the station's two-year return-period level, a threshold computed exclusively from the annual maxima of the training years to prevent data leakage. To ensure rigorous evaluation, we employed an expanding-window walk-forward backtest where the model for a given year Y is trained only on data from 1997 to Y-1, with test years spanning 2021 to 2026. This procedure allows the model to be run day by day through the recent monsoon seasons, providing a realistic assessment of operational performance in a data-sparse environment.

For impact modeling, we constructed a logistic regression model that predicts district-level flood impacts reported by the Assam State Disaster Management Authority (ASDMA). The unit of analysis is the district-month during the monsoon season (May to October) for the period 2021 to 2026, covering thirty-three districts. The target variable is whether the reported affected population exceeds a minimum threshold, MIN\_AFFECTED. Features include river forecast signals from the walk-forward backtest, district rainfall, and historical flood records, all of which are available before the month's impacts are known. We validated these forecasts against real flood impacts using ASDMA district reports compiled by CivicDataLab's IDS-DRR project, which provide ground truth for the population affected per district and month from April 2021 to September 2026. A forecast is considered a warning if, on any day of the month, the three-day forecast for the district's river station crosses its alert cut-off.

For satellite flood mapping, we implemented a U-Net architecture with a ResNet-18 encoder, trained from scratch on Sentinel-1 radar imagery. The training data consists of hand-labelled chips from the Sen1Floods11 dataset, which includes 446 scenes from eleven flood events, including those in Assam. We compared our model against Otsu thresholding on the VH polarization, a classic operational baseline. Additionally, we evaluated the forecasting system against several baselines, including a river-only forecast without rainfall, a heavy local rain rule based on the top ten percent of three-day totals, and a simple calendar baseline that flags every June and July. For daily validation, we compared the river-only model, the river plus rainfall model, and the rain-only reference against ASDMA's daily district reports for the 2022 monsoon. A district is classified as flooded on a given day when ASDMA reports at least MIN\_AFFECTED people affected, allowing us to assess the timing of flood onset relative to the forecast warnings.

\section{Evaluation Design}
We validated the forecasting pipeline using an expanding-window walk-forward backtest to prevent data leakage, ensuring that for any given test year Y, the model was trained exclusively on data from 1997 to Y-1. The evaluation period comprised the monsoon seasons from 2021 to 2026. For river flood warnings, a flood event was defined as river flow exceeding the station's two-year return-period level, a threshold computed strictly from the annual maxima of the training years to maintain temporal integrity. We assessed performance using the Critical Success Index (CSI) and compared our gradient-boosted models against several baselines. These included a river-only forecast, a heavy local rain rule based on the top 10\% of three-day rainfall totals, and a simple calendar baseline that flagged every June and July. Additionally, we evaluated the models against a ceiling baseline representing the scenario where the river was actually above the warning level, serving as an observational upper bound rather than a predictive forecast.

For impact modeling, the unit of evaluation was the district-month during the monsoon (May–October) across 33 districts. A district-month was classified as flooded if the Assam State Disaster Management Authority (ASDMA) reported at least a minimum number of affected people. We validated these forecasts against real-world impacts using ASDMA district reports compiled from April 2021 to September 2026. A forecast was considered a successful warning if, on any day within the month, the three-day river forecast for the district's station crossed its alert cut-off. We also conducted a day-by-day validation against ASDMA's daily district reports for the 2022 monsoon to assess early-warning performance. In this daily check, a district was deemed flooded on a specific day if ASDMA reported the minimum threshold of affected individuals. We compared the river-only model, the combined river and rainfall model, and the simple rain-only reference to determine the incremental value of multi-source data integration.

\section{Results}
We evaluated the gradient-boosted river flood forecasting model using an expanding-window walk-forward backtest, where the model was trained on data from 1997 to Y-1 to forecast year Y, with test years spanning 2021 to 2026. The flood definition was set as flow exceeding the station's 2-year return-period level, computed from the annual maxima of the training years only. All reported metrics are out-of-sample, using the Critical Success Index (CSI), defined as hits divided by the sum of hits, misses, and false alarms. For the 1-day lead time, the model achieved a CSI of 0.82, catching 90\% of flood days with a 10\% false alarm rate. This performance slightly exceeded the river-flow-only baseline, which recorded a CSI of 0.81, and significantly outperformed the persistence baseline of assuming tomorrow equals today, which yielded a CSI of 0.73.

As the lead time increased, the forecasting skill decreased, reflecting the inherent uncertainty in long-range hydrological predictions. At a 2-day lead time, the model achieved a CSI of 0.72, with 80\% of flood days caught and a 12\% false alarm rate. The river-flow-only baseline performed slightly better at this lead time with a CSI of 0.73, while the persistence baseline dropped to 0.65. For a 3-day lead time, the combined model and the river-flow-only baseline both achieved a CSI of 0.64, catching 73\% of flood days with a 16\% false alarm rate. The persistence baseline continued to degrade, reaching a CSI of 0.59. At a 5-day lead time, the model recorded a CSI of 0.51, with 65\% of flood days caught and a 31\% false alarm rate, matching the river-flow-only baseline exactly. The persistence baseline achieved a CSI of 0.49. Finally, at a 7-day lead time, the model achieved a CSI of 0.43, catching 59\% of flood days with a 39\% false alarm rate, while the river-flow-only baseline scored 0.45 and the persistence baseline scored 0.42.

We assessed the early warning capability of the model by validating forecasts against ASDMA's daily district reports for the 2022 monsoon, specifically examining 45 district flood onsets defined as the first day with at least 1,000 people affected after a quiet week. The river-plus-rainfall model warned of 51\% of onsets one day in advance, 42\% three days in advance, and 18\% five days in advance, with a false alarm rate of 0.9\% on quiet days. In contrast, the river-flow-only model warned of 38\% of onsets one day in advance, 11\% three days in advance, and 2\% five days in advance, also with a 0.9\% false alarm rate. A simple rule based on heavy local rain, defined as the top 10\% of 3-day totals, warned of 78\% of onsets one day in advance, 71\% three days in advance, and 58\% five days in advance, but with a higher false alarm rate of 1\%. The observation that the river was already above 85\% of the warning level warned of 42\% of onsets one day in advance, 13\% three days in advance, and 4\% five days in advance.

We evaluated the impact model's ability to predict district-level reported flood impacts using ASDMA district reports from May to October 2021 to 2026, with the unit of analysis being the district-month. The impact model, which combined river forecasts, local rainfall, and district history, achieved a CSI of 0.41, warning of 60\% of floods with a 44\% false alarm rate. This outperformed the river forecast alone, which achieved a CSI of 0.32 with 42\% of floods warned and a 44\% false alarm rate. The climatology baseline, which flagged every June and July, achieved a CSI of 0.32 with 55\% of floods warned and a 57\% false alarm rate. Using district flood history alone resulted in a CSI of 0.30, warning of 42\% of floods with a 47\% false alarm rate. The observation that the river was actually above the warning level, serving as a ceiling for forecast performance, achieved a CSI of 0.30 with 37\% of floods warned and a 41\% false alarm rate.

We evaluated the U-Net flood water segmentation model on the Sen1Floods11 dataset, which includes 446 hand-labelled scenes from 11 flood events. On the full test set covering 10 countries, the U-Net achieved a Water IoU of 0.65 and a pixel accuracy of 95\%, substantially outperforming the Otsu thresholding baseline, which achieved a Water IoU of 0.23 and a pixel accuracy of 70\%. For Assam-specific scenes, the U-Net achieved a Water IoU of 0.72. To assess generalization to unseen events, we tested the model on Bolivia scenes, which were not included in the training data, where the U-Net achieved a Water IoU of 0.63 compared to the Otsu baseline's 0.40. These results indicate that the deep learning approach provides robust flood extent mapping across diverse geographical contexts, significantly improving upon the operational Otsu thresholding method.

\section{Discussion and Limitations}
The results indicate that integrating multi-scale rainfall features with river discharge data enhances short-term flood warning performance in the Brahmaputra basin. Our gradient-boosted model achieved a Critical Success Index (CSI) of 0.82 for one-day forecasts, outperforming the river-only baseline of 0.81 and the persistence baseline of 0.73. This aligns with recent literature suggesting that hybrid approaches and ensemble methods improve flood prediction accuracy ~\cite{mosavi2018c1,kan2020c7}. However, performance degrades with lead time; at seven days, the CSI drops to 0.43, which is lower than the river-only baseline of 0.45. This suggests that while rainfall data adds value for immediate warnings, its predictive power diminishes for longer horizons, consistent with findings that deep learning models may offer better stability across different lead times ~\cite{hu2018c19}. The day-by-day early warning analysis against ASDMA 2022 reports shows that the combined model detected 51\% of flood events one day in advance, compared to 38\% for the river-only model.

A significant limitation is the reliance on reanalysis data rather than in-situ measurements. River flow inputs are derived from GloFAS, and rainfall is sampled at only 15 catchment points, which may limit the capture of highly localized downpours. This data sparsity is a common challenge in monsoon-affected regions, where machine learning models often face limitations in generalization ~\cite{bayat2020c11}. Furthermore, the impact model, which predicts district-level affected populations, relies on ASDMA reports available only from 2021 to 2026. Under a strict leave-one-year-out validation, the impact model achieved a CSI of 0.41, improving upon the river-only baseline of 0.32. However, the model's performance in 2022 was limited to a CSI of 0.13 due to the lack of historical data, highlighting the dependency on sufficient historical records for impact prediction. This contrasts with studies using LSTM networks that demonstrate capabilities in predicting flood severity with real-time data ~\cite{mourot2022c14}.

For satellite-based flood mapping, the U-Net model achieved a Water Intersection over Union (IoU) of 0.72 on Assam scenes, substantially outperforming the Otsu threshold baseline. On the broader Sen1Floods11 test set, the U-Net attained a Water IoU of 0.65, compared to 0.23 for the Otsu method. These results support the trend that deep learning models produce better accuracy than traditional machine learning or thresholding approaches for inundation modeling ~\cite{karim2023c3}. However, the model's performance on unseen events, such as the Bolivia scenes, was lower at 0.63, indicating potential issues with generalization to different geographic contexts. This aligns with the observation that while large-sample datasets can improve model robustness, local calibration remains important ~\cite{kratzert2019c16}. The pixel accuracy of 95\% for the U-Net on the test set further confirms its utility for operational mapping, though the reliance on Sentinel-1 data may introduce uncertainties in cloud-covered or complex terrain scenarios.

\section{Conclusion}
We presented a gradient-boosted framework for short-term river and rainfall forecasting, coupled with impact modeling and Sentinel-1 flood mapping for the Brahmaputra basin in Assam. Walk-forward backtests from 2021 to 2026 demonstrated that the river-level exceedance model achieved a Critical Success Index of 0.82 for one-day lead times, outperforming simple persistence and river-only baselines. The impact model, which integrates river forecasts with district rainfall and historical records, improved district-month impact CSI to 0.41 compared to river-only and historical baselines. Additionally, the U-Net segmentation model attained a Water IoU of 0.72 on Assam scenes, substantially exceeding the Otsu threshold baseline of 0.23 observed in the broader test set. These results indicate that combining hydrological forecasting with satellite-derived extent mapping provides a robust operational tool for monsoon flood management.

Several limitations constrain the current system. River flow inputs rely on GloFAS reanalysis rather than in-situ gauge measurements, and rainfall is sampled at only 15 points, which may limit the capture of highly localized downpours. Furthermore, ASDMA impact data are limited to the 2021–2026 period, with daily detail available for only one year, restricting the temporal depth of impact model training. Future work should incorporate higher-resolution in-situ hydrological data and expand the temporal coverage of district-level impact reports to improve model generalizability. Additionally, extending the satellite mapping pipeline to include multi-temporal analysis could enhance the detection of flood dynamics in data-sparse regions.

\bibliographystyle{IEEEtran}
\bibliography{refs}
\end{document}